Este experimento consistió en la determinación del coeficiente de expansión lineal para varillas huecas de cobre y acero mediante un montaje que permitió medir cambios longitudinales en función del aumento de temperatura. El fundamento teórico se basa en que el incremento de la energía cinética de las moléculas provoca un aumento en su distancia promedio de separación, fenómeno que se manifiesta como una dilatación proporcional al cambio térmico. Para la recolección de datos, se utilizó una bomba de agua con calentador para circular líquido a través de las varillas y un calibrador de reloj para registrar el desplazamiento mecánico con alta precisión. Tras realizar las regresiones lineales sobre una longitud inicial establecida, se obtuvieron los coeficientes experimentales específicos para cada metal. Estos resultados presentan una alta confiabilidad, respaldada por coeficientes de correlación cercanos a la unidad para ambos materiales. En conclusión, se validó la proporcionalidad directa entre la temperatura y la elongación física de los materiales, confirmando que la metodología y la precisión del equipo utilizado permitieron una caracterización térmica efectiva de los metales analizados en el laboratorio.
